\documentclass[11pt,a4paper]{article}
\usepackage[T1]{fontenc}
\usepackage[utf8]{inputenc}
\usepackage{mathtools,amsthm}
\usepackage{newtxtext,newtxmath}
\usepackage[a4paper,margin=25mm,headheight=15pt]{geometry}
\usepackage{microtype}
\usepackage{booktabs,array,longtable}
\usepackage[dvipsnames]{xcolor}
\usepackage{enumitem}
\usepackage{fancyhdr}
\usepackage{titlesec}
\usepackage{listings}
\usepackage{hyperref}
\definecolor{inkblue}{HTML}{173B57}
\definecolor{muted}{HTML}{526576}
\definecolor{lightblue}{HTML}{F1F5F8}
\hypersetup{colorlinks=true,linkcolor=inkblue,citecolor=inkblue,urlcolor=inkblue,
 pdftitle={Shifted Catalan Hankel Polynomials: Exact Coefficients and Sharp Recurrences},
 pdfauthor={OpenAI ChatGPT},pdfsubject={A proof of a denominator conjecture and an indexing correction}}
\titleformat{\section}{\large\bfseries\color{inkblue}}{\thesection}{0.7em}{}
\titleformat{\subsection}{\normalsize\bfseries\color{inkblue}}{\thesubsection}{0.7em}{}
\pagestyle{fancy}\fancyhf{}
\fancyhead[L]{\small\color{muted}Shifted Catalan Hankel polynomials}
\fancyhead[R]{\small\color{muted}Research note \textbar\ September 2026}
\fancyfoot[C]{\thepage}
\renewcommand{\headrulewidth}{0.3pt}
\setlength{\parskip}{4pt}
\setlength{\parindent}{1.2em}
\setlength{\emergencystretch}{2em}
\allowdisplaybreaks[2]
\numberwithin{equation}{section}
\newtheorem{theorem}{Theorem}[section]
\newtheorem{lemma}[theorem]{Lemma}
\newtheorem{proposition}[theorem]{Proposition}
\newtheorem{corollary}[theorem]{Corollary}
\theoremstyle{definition}
\newtheorem{definition}[theorem]{Definition}
\newtheorem{example}[theorem]{Example}
\theoremstyle{remark}
\newtheorem{remark}[theorem]{Remark}
\newcommand{\NN}{\mathbb N_0}
\newcommand{\CC}{\mathbb C}
\newcommand{\ZZ}{\mathbb Z}
\newcommand{\QQ}{\mathbb Q}
\newcommand{\D}[2]{D_{#1}^{(#2)}}
\newcommand{\HH}[2]{H_{#1}^{(#2)}}
\newcommand{\dd}{\mathrm d}
\newcommand{\lc}{\operatorname{lc}}
\newcommand{\sgn}{\operatorname{sgn}}
\newcommand{\hyp}[4]{{}_2F_1\!\left(\begin{matrix}#1,#2\\#3\end{matrix};#4\right)}
\lstset{basicstyle=\small\ttfamily,breaklines=true,columns=fullflexible,
 frame=single,rulecolor=\color{muted},backgroundcolor=\color{lightblue},
 keepspaces=true,showstringspaces=false}

\begin{document}
\thispagestyle{empty}
\begin{center}
{\small\sffamily\bfseries\color{muted}EXACT ENUMERATION \quad / \quad GENERATING FUNCTIONS \quad / \quad ASYMPTOTICS}\par
\vspace{10mm}
{\LARGE\bfseries\color{inkblue}Shifted Catalan Hankel Polynomials}\par
\vspace{3mm}
{\Large Exact coefficients, sharp recurrences,\par and a denominator conjecture}\par
\vspace{7mm}
{A research note prepared with OpenAI ChatGPT}\par
{19 September 2026}
\end{center}
\vspace{5mm}
\begin{abstract}
For the Catalan numbers $C_n$, consider the size-$N$ determinant
$D_N^{(m)}(a,b)=\det(aC_{i+j+m}+bC_{i+j+m+1})_{0\leq i,j<N}$.
We give a self-contained proof of the denominator-exponent conjecture stated
in Paul Barry's 2020 note on this family.  With $d=\binom m2$, the generating
function is $P_m(t)/(1-(a+2b)t+b^2t^2)^{d+1}$.  The denominator is minimal
whenever $b\ne0$, $a\ne0$, and $a\ne-4b$.  We determine every exceptional
case, prove that $\deg P_m=(m-1)^2$ for $b\ne0$, and establish a reciprocal
symmetry of the numerator.  An explicit positive-integer coefficient formula
also identifies and repairs a one-unit shift in the source's separate
Conjecture~2.  The essential degree reduction is a Vandermonde cancellation
in confluent Christoffel minors.  It yields an exact sum of two exponentials
with polynomial coefficients, a finite full asymptotic expansion, and an
explicit first correction.  Independent integer and rational computations,
including direct determinants and polynomial-identity certificates at fixed
sizes, accompany the proofs.
\end{abstract}
\vspace{3mm}
\noindent\colorbox{lightblue}{\parbox{\dimexpr\textwidth-2\fboxsep\relax}{%
\textbf{Main result.} For all nonexceptional complex parameters,
\[
\boxed{\ \text{minimal constant-coefficient recurrence order}
       =m(m-1)+2.\ }
\]
The proof is uniform in the determinant size and the Catalan shift.
}}\par
\vspace{5mm}
\noindent\textbf{Status and scope.}
This is an unrefereed, AI-assisted mathematical manuscript, not a formalized
proof.  Its theorems are supported by explicit arguments, not inferred from
finite tests.  The target is a conjecture in a specified source version;
no claim of first discovery or exhaustive verification of present-day
open status is made.  Classical Christoffel formulas and the known general
rationality theory are credited, rather than presented as new.
\vfill
\noindent\textbf{Reproducibility.} The accompanying archive contains this
\LaTeX\ source, the PDF, a standard-library-only Python implementation,
an independently organized verification suite, and exact machine-readable
data.  No external mathematical software is needed to rerun the tests.

\clearpage
\begingroup
\small
\setlength{\parskip}{0pt}
\tableofcontents
\endgroup
\clearpage

\section{The problem and the contribution}\label{sec:problem}

Let
\begin{equation}\label{eq:catalan}
 C_n=\frac{1}{n+1}\binom{2n}{n}\qquad(n\in\NN)
\end{equation}
be the Catalan numbers.  We use determinant size, rather than the last row
index, throughout:
\begin{equation}\label{eq:D}
 \D{N}{m}(a,b)=\det_{0\leq i,j<N}
       \bigl(aC_{i+j+m}+bC_{i+j+m+1}\bigr),
 \qquad \D{0}{m}=1.
\end{equation}
Here $m,N\in\NN$, and the parameters may be indeterminates or complex
numbers.  Define the ordinary generating function
\begin{equation}\label{eq:F}
 F_m(t;a,b)=\sum_{N\geq0}\D{N}{m}(a,b)t^N.
\end{equation}
All generating-function identities below first make sense as identities of
formal power series.  Assertions about poles concern the associated
rational functions.

Barry's preprint~\cite{Barry2020}, version~1, explicitly conjectures on
page~4 that the exponent of the quadratic generating-function denominator
is $\binom m2+1$.  The same page contains a more complicated coefficient
conjecture.  These are distinct statements: the denominator claim is
correct in the uniform sense proved here, whereas the printed coefficient
formula uses a shift one smaller than its accompanying determinant.
Appendix~\ref{app:barry} makes the translation precise.

There is substantial prior theory.  Dougherty, French, Saderholm, and
Qian~\cite{Dougherty2011} studied recurrences for linear combinations of
Catalan numbers.  Krattenthaler~\cite{Krattenthaler2023} gives a detailed
account of the Christoffel-type determinant identity, confluence, and
constant-coefficient recurrences; his Section~8 also records Elouafi's
prior proof of general Catalan recurrence assertions.  For a multiplier
of degree $r$, that general construction supplies an order at most $2^r$.
Consequently, rationality by itself is \emph{not} the research claim here.
For the special multiplier $x^m(a+bx)$, the objective is the exact,
much smaller order, together with its sharp parameter classification.

The article separates two useful calculations.  Jacobi orthogonality
immediately exposes the individual coefficients.  A confluent determinant,
on the other hand, exposes the dependence on $N$ and explains the precise
recurrence order.  Their independent implementations provide a useful
check against indexing and normalization errors.

\subsection{A notation guide}
We shall repeatedly use
\begin{equation}\label{eq:notation}
 d=\binom m2,\qquad e=d+1,\qquad K=(m-1)^2,\qquad
 q(t)=1-(a+2b)t+b^2t^2.
\end{equation}
The shifted pure Catalan determinant and its leading coefficient are
\begin{align}
 \HH{N}{m}&=\det_{0\leq i,j<N}(C_{i+j+m})
 =\prod_{1\leq j\leq i<m}\frac{2N+i+j}{i+j},\label{eq:Hproduct}\\
 h_m&=\prod_{1\leq j\leq i<m}\frac{2}{i+j}.
 \label{eq:h}
\end{align}
Every empty product and empty determinant equals~$1$.  In particular,
$H_N^{(0)}=H_N^{(1)}=1$ and $h_0=h_1=1$.  We prove
\eqref{eq:Hproduct} rather than assuming it.  As a polynomial in $N$,
$H_N^{(m)}$ has degree $d$ and leading coefficient $h_m$.
The notation $(u)_k=u(u+1)\cdots(u+k-1)$ denotes a rising factorial,
with $(u)_0=1$.

\section{Statements of the main theorems}\label{sec:main}

\begin{theorem}[Explicit coefficients]\label{thm:coeff}
For every $N,m\in\NN$,
\begin{equation}\label{eq:coeff}
 \D{N}{m}(a,b)=\sum_{k=0}^N c_{N,m,k}a^k b^{N-k},\qquad
 c_{N,m,k}=\HH{N}{m+1}\binom Nk
       \frac{(N+m+1)_k}{4^k(m+\tfrac12)_k}.
\end{equation}
All the coefficients $c_{N,m,k}$ are positive integers.  Equivalently,
\begin{equation}\label{eq:coeff-factorial}
 c_{N,m,k}=\HH{N}{m+1}\binom Nk
 \frac{(N+m+k)!(2m)!(m+k)!}
 {(N+m)!(2m+2k)!m!}.
\end{equation}
They can be generated without evaluating any determinant:
\begin{equation}\label{eq:coeff-update}
 c_{N,m,0}=\HH{N}{m+1},\qquad
 \frac{c_{N,m,k+1}}{c_{N,m,k}}=
 \frac{(N-k)(N+m+1+k)}{2(k+1)(2m+2k+1)}.
\end{equation}
\end{theorem}

\begin{theorem}[Uniform rational form and sharp denominator]\label{thm:main}
For every $m\in\NN$ there is a polynomial $P_m\in\ZZ[a,b,t]$ such that
\begin{equation}\label{eq:mainGF}
 F_m(t;a,b)=\frac{P_m(t;a,b)}{q(t)^{d+1}}.
\end{equation}
It satisfies $P_m(0;a,b)=1$ and $\deg_t P_m\leq K$.
If $b\ne0$, then $\deg_t P_m=K$ and
\begin{equation}\label{eq:reciprocity}
 P_m(t;a,b)=s_m(bt)^K P_m\!\left(\frac{1}{b^2t};a,b\right),
 \qquad s_m=(-1)^{(m-1)(m-2)/2}.
\end{equation}
In particular, $[t^K]P_m=s_m b^K$.
If $b\ne0$, $a\ne0$, and $a\ne-4b$, the representation
\eqref{eq:mainGF} is reduced over $\CC[t]$.  The sequence has minimal
constant-coefficient recurrence characteristic polynomial
\begin{equation}\label{eq:charpoly}
 \bigl(X^2-(a+2b)X+b^2\bigr)^{d+1},
\end{equation}
and hence minimal order $2d+2=m(m-1)+2$.
\end{theorem}

Here ``minimal'' refers even to recurrences allowed to start after a
finite initial segment.  Thus minimality is not an artifact of the
empty-determinant convention.  The polynomial numerator in
\eqref{eq:mainGF} remains the \emph{uniform, possibly unreduced} numerator
at exceptional parameters.

\begin{theorem}[Every exceptional denominator]\label{thm:exceptional}
For fixed complex $a,b$, the reduced denominator of $F_m$, normalized to
have constant term $1$, is exactly the following:
\begin{center}
\renewcommand{\arraystretch}{1.35}
\begin{tabular}{@{}ll@{}}
\toprule
Parameter condition & Reduced denominator\\
\midrule
$a=b=0$ & $1$\\
$b=0$, $a\ne0$ & $(1-at)^{\binom m2+1}$\\
$a=0$, $b\ne0$ & $(1-bt)^{\binom{m+1}{2}+1}$\\
$a=-4b$, $b\ne0$ & $(1+bt)^{\binom m2+2}$\\
$b\ne0$, $a(a+4b)\ne0$ & $q(t)^{\binom m2+1}$\\
\bottomrule
\end{tabular}
\end{center}
\end{theorem}

\begin{theorem}[Finite asymptotic expansion]\label{thm:asymp}
Let $a,b>0$ be fixed and put
\begin{equation}\label{eq:positive-roots}
 r=\frac{a/b+2+\sqrt{(a/b)(a/b+4)}}2>1,\qquad
 \lambda_+=br,\quad \lambda_-=b/r.
\end{equation}
There are explicit degree-$d$ polynomials $B_+(N),B_-(N)$ such that
\begin{equation}\label{eq:two-exp-statement}
 \D{N}{m}(a,b)=B_+(N)\lambda_+^N+B_-(N)\lambda_-^N
 \qquad(N\geq0).
\end{equation}
In particular,
\begin{align}
 \D{N}{m}(a,b)
 &=A_m N^d\lambda_+^N
 \left(1+\frac{d\left(\frac m2-\frac1{r-1}\right)}{N}
       +O(N^{-2})+O(r^{-2N})\right),\label{eq:asymp}\\
 A_m&=h_m\frac{r^{m+1}}{(r-1)^m(r+1)}.\label{eq:A}
\end{align}
The entire algebraic expansion terminates after $N^{-d}$; the second
exponential is also multiplied by a polynomial of degree~$d$.
For $d=0$, there are no algebraic correction terms at all.
\end{theorem}

\subsection{The scale of the improvement}
For a degree-$(m+1)$ multiplier the general $2^{m+1}$ construction need
not be minimal.  The theorem determines the actual order in this family.
The first values are:
\begin{center}
\renewcommand{\arraystretch}{1.12}
\begin{tabular}{@{}rrrrr@{}}
\toprule
$m$ & Exponent $e$ & Exact order $2e$ & Numerator degree $K$ & $2^{m+1}$\\
\midrule
0 & 1 & 2 & 1 & 2\\
1 & 1 & 2 & 0 & 4\\
2 & 2 & 4 & 1 & 8\\
3 & 4 & 8 & 4 & 16\\
4 & 7 & 14 & 9 & 32\\
5 & 11 & 22 & 16 & 64\\
6 & 16 & 32 & 25 & 128\\
8 & 29 & 58 & 49 & 512\\
10 & 46 & 92 & 81 & 2048\\
\bottomrule
\end{tabular}
\end{center}
The reduction is not obtained by fitting a shorter recurrence to a long
initial segment.  It comes from a cancellation of polynomial degree in a
family of minors, proved in Section~\ref{sec:degree}.

\section{Moments, orthogonality, and the coefficient formula}\label{sec:coeffproof}

\subsection{Catalan moments and a normalized Jacobi polynomial}
For $m\in\NN$ define a positive measure on $(0,4)$ by
\begin{equation}\label{eq:measure}
 \dd\mu_m(x)=\frac1{2\pi}x^{m-1/2}(4-x)^{1/2}\,\dd x.
\end{equation}
The beta integral, after $x=4t$, gives
\begin{equation}\label{eq:moments}
 \int_0^4 x^n\,\dd\mu_m(x)
 =\frac{4^{n+m+1}}{2\pi}B(n+m+\tfrac12,\tfrac32)
 =C_{n+m}.
\end{equation}
Consequently every moment matrix is positive definite, so its determinant
is nonzero and its monic orthogonal polynomials are well-defined.

\begin{lemma}[The monic polynomial]\label{lem:Jacobi}
The degree-$N$ monic orthogonal polynomial for $\mu_m$ is
\begin{equation}\label{eq:Jacobi}
 p_N^{(m)}(x)=(-1)^N
 \frac{4^N(m+\tfrac12)_N}{(N+m+1)_N}
 \hyp{-N}{N+m+1}{m+\tfrac12}{x/4}.
\end{equation}
The hypergeometric series terminates after $N+1$ terms.
\end{lemma}
\begin{proof}
Write $\alpha=m-\tfrac12$, $\beta=\tfrac12$, so both are greater than $-1$.
A Rodrigues representation on $0<t<1$ is
\begin{equation}\label{eq:Rodrigues}
 \frac{t^{-\alpha}(1-t)^{-\beta}}{(\alpha+1)_N}
 \frac{\dd^N}{\dd t^N}
       \left(t^{N+\alpha}(1-t)^{N+\beta}\right)
 =\hyp{-N}{N+\alpha+\beta+1}{\alpha+1}{t}.
\end{equation}
Leibniz's rule gives the polynomial on the right, including its
normalization at $t=0$.  Multiplying the left side by
$t^j t^\alpha(1-t)^\beta$ and integrating by parts $N$ times yields
zero for $0\leq j<N$.  Every boundary term vanishes: after at most
$N-1$ derivatives, the endpoint exponents that can occur are at least
$\alpha+1>0$ and $\beta+1>0$.  The final derivative of $t^j$ is zero.

The coefficient of $t^N$ on the right is
$(-1)^N(N+\alpha+\beta+1)_N/(\alpha+1)_N$.
Replacing $t$ by $x/4$ and normalizing this leading coefficient gives
\eqref{eq:Jacobi}.  This also identifies the polynomial with the standard
Jacobi family; see the conventional hypergeometric representation in
DLMF~\cite{DLMF}, Eq.~18.5.7.
\end{proof}

\subsection{A characteristic-polynomial determinant identity}
The following elementary identity will be used twice.

\begin{lemma}\label{lem:characteristic}
Let $\nu_j$ be the moments of a positive measure, let
$H_N=\det(\nu_{i+j})_{0\leq i,j<N}$, and let $p_N$ be its monic orthogonal
polynomial.  Then
\begin{equation}\label{eq:characteristic}
 \det_{0\leq i,j<N}(\nu_{i+j+1}-z\nu_{i+j})
   =(-1)^N H_Np_N(z).
\end{equation}
\end{lemma}
\begin{proof}
The monic polynomial is represented by
\[
 p_N(z)=\frac1{H_N}
 \det\begin{pmatrix}
 \nu_0&\nu_1&\cdots&\nu_N\\
 \nu_1&\nu_2&\cdots&\nu_{N+1}\\
 \vdots&\vdots&&\vdots\\
 \nu_{N-1}&\nu_N&\cdots&\nu_{2N-1}\\
 1&z&\cdots&z^N
 \end{pmatrix}.
\]
Its leading coefficient is one, and integrating it against any power
$z^j$, $j<N$, creates a repeated row.  Thus it is indeed $p_N$.
Starting with the last column and moving left, replace
$\mathrm{col}_j$ by $\mathrm{col}_j-z\,\mathrm{col}_{j-1}$.
The last row becomes $(1,0,\ldots,0)$.  Expansion along it gives
\eqref{eq:characteristic}, with sign $(-1)^N$.
\end{proof}

Applying the lemma to $\mu_m$ and $z=-a/b$ gives, for $b\ne0$,
\begin{equation}\label{eq:characteristicD}
 \D{N}{m}(a,b)=(-b)^N\HH{N}{m}\,p_N^{(m)}(-a/b).
\end{equation}

\subsection{A proof of the shifted Catalan product}
For $m=0$, \eqref{eq:Jacobi} reduces to
\begin{equation}\label{eq:p0}
 p_n^{(0)}(x)=\sum_{j=0}^n(-1)^{n-j}\binom{n+j}{2j}x^j.
\end{equation}
Its initial values and recurrence are
\begin{equation}\label{eq:p0rec}
 p_0^{(0)}=1,\quad p_1^{(0)}=x-1,\quad
 p_{n+2}^{(0)}=(x-2)p_{n+1}^{(0)}-p_n^{(0)}.
\end{equation}
These follow directly by comparing the binomial coefficients in
\eqref{eq:p0}.  If $\eta_n$ denotes the squared norm of $p_n^{(0)}$, the
coefficient $1$ of $p_{n-1}^{(0)}$ in $xp_n^{(0)}$ implies
$\eta_n=\eta_{n-1}$.  Indeed, the two ways of writing
$\langle xp_n^{(0)},p_{n-1}^{(0)}\rangle$
give $\eta_{n-1}$ and $\eta_n$, respectively.  Since
$\eta_0=C_0=1$, all these norms equal~$1$.  Changing from monomials to
monic orthogonal polynomials is unitriangular, and therefore
$H_N^{(0)}=\prod_{n=0}^{N-1}\eta_n=1$.

Next set $z=0$ in \eqref{eq:characteristic} and use
\eqref{eq:Jacobi}.  We obtain
\begin{equation}\label{eq:Hratio}
 \frac{\HH{N}{m+1}}{\HH{N}{m}}
 =(-1)^Np_N^{(m)}(0)
 =\frac{4^N(m+\tfrac12)_N}{(N+m+1)_N}
 =\prod_{j=1}^{m}\frac{2N+m+j}{m+j}.
\end{equation}
For completeness, the middle expression equals
\[
 \frac{(2m+2N)!\,m!\,(N+m)!}
      {(m+N)!\,(2m)!\,(2N+m)!}
 =\frac{(2N+2m)!\,m!}{(2m)!\,(2N+m)!},
\]
which is the last product.  Induction on $m$, beginning with
$H_N^{(0)}=1$, proves \eqref{eq:Hproduct}.

\subsection{Completion of the coefficient proof}
Insert \eqref{eq:Jacobi} into \eqref{eq:characteristicD}.
The constant factor is $H_N^{(m+1)}$ by \eqref{eq:Hratio}.
In each summand, $(-N)_k(-1)^k/k!=\binom Nk$.  This gives exactly
\eqref{eq:coeff}.  The formula is a polynomial identity for $b\ne0$,
so it holds for $b=0$ as well.  The duplication identity for a rising
half-integer factorial, or direct cancellation, gives
\eqref{eq:coeff-factorial}.  Taking adjacent coefficient ratios gives
\eqref{eq:coeff-update}.

All factors in \eqref{eq:coeff} are positive rational numbers.  The
left side is a determinant with entries in $\ZZ[a,b]$, so its
coefficients are integers.  This proves positivity and integrality
simultaneously and completes Theorem~\ref{thm:coeff}.
In particular,
\begin{equation}\label{eq:extremalcoeff}
 c_{N,m,0}=\HH{N}{m+1},\qquad c_{N,m,N}=\HH{N}{m}.
\end{equation}

\subsection{Real roots and coefficient inequalities}
\begin{corollary}\label{cor:roots}
For $N\geq1$, the polynomial $a\mapsto D_N^{(m)}(a,1)$ has exactly $N$
simple real roots, all in $(-4,0)$.  For $N\geq2$, its normalized
coefficients $u_k=c_{N,m,k}/\binom Nk$ satisfy
\begin{equation}\label{eq:ultra}
 u_k^2>u_{k-1}u_{k+1}\qquad(1\leq k<N).
\end{equation}
\end{corollary}
\begin{proof}
A degree-$N$ orthogonal polynomial for a strictly positive density on
$(0,4)$ has $N$ simple roots in that interval.  One elementary proof
is to multiply it by the product of its distinct sign-changing roots
in $(0,4)$.  If there were fewer than $N$ such roots, that product
would have degree below $N$, while the integral of its product with
the orthogonal polynomial would have constant nonzero sign.  This
contradicts orthogonality.  Equation~\eqref{eq:characteristicD} now
negates the roots.

Alternatively, the coefficient inequality follows without any root
theory.  Theorem~\ref{thm:coeff} gives
\[
 \frac{u_{k+1}}{u_k}
 =\frac{N+m+1+k}{4(m+\tfrac12+k)}.
\]
This is strictly decreasing in $k$, because the numerator's constant
term exceeds the denominator's unscaled constant term by $N+1/2$.
The decrease is equivalent to \eqref{eq:ultra}.
\end{proof}

\section{A fixed-size confluent determinant}\label{sec:confluent}

The explicit coefficient formula grows in length with $N$.  To understand
the recurrence order, we need a representation whose determinant size
depends only on $m$.
Define
\begin{equation}\label{eq:R}
 R_n(y)=(-1)^np_n^{(0)}(-y)
       =\sum_{i=0}^n\binom{n+i}{2i}y^i.
\end{equation}
Thus
\begin{equation}\label{eq:Rrec}
 R_0(y)=1,\quad R_1(y)=y+1,\quad
 R_{n+2}(y)=(y+2)R_{n+1}(y)-R_n(y).
\end{equation}
Extend $R_n$ to all integer indices using this recurrence.  Since
$R_{-1}=R_0=1$, it also satisfies the reflection identity
\begin{equation}\label{eq:Rreflection}
 R_{-n-1}(y)=R_n(y)\qquad(n\in\ZZ).
\end{equation}
The reversed form of the polynomials is convenient because it eliminates
unnecessary signs from Christoffel's identity.

\begin{lemma}[Christoffel determinant in the present normalization]
\label{lem:Christoffel}
For pairwise distinct $y_1,\ldots,y_r$,
\begin{equation}\label{eq:Christoffel}
 \det_{0\leq i,j<N}\int_0^4 x^{i+j}
        \prod_{\ell=1}^r(x+y_\ell)\,\dd\mu_0(x)
 =\frac{\det\bigl[R_{N+j}(y_i)\bigr]_{\substack{1\leq i\leq r\\0\leq j<r}}}
        {\prod_{1\leq i<j\leq r}(y_j-y_i)}.
\end{equation}
The quotient extends polynomially when some $y_i$ coincide.
\end{lemma}
\begin{proof}
Here is a proof that also verifies the normalization.  Write
$\Delta(x_1,\ldots,x_N)=\det[x_i^{j-1}]_{1\leq i,j\leq N}$.
Expanding two determinants and integrating termwise gives the identity
\[
 \frac1{N!}\int\cdots\int
 \det[f_i(x_j)]\det[g_i(x_j)]\prod_j\dd\mu_0(x_j)
 =\det\left[\int f_i(x)g_j(x)\,\dd\mu_0(x)\right].
\]
No analytic interchange difficulty occurs here: the sums are finite and
all polynomials are integrable on a bounded interval.

For external variables $z_1,\ldots,z_r$ we have
\[
 \Delta(x,z)=\Delta(x)\Delta(z)
            \prod_{i=1}^N\prod_{\ell=1}^r(z_\ell-x_i).
\]
In the determinant defining $\Delta(x,z)$, replace monomials by the
monic polynomials $p_j^{(0)}$; the determinant is unchanged.  Expand
along its first $N$ rows and integrate against $\Delta(x)$.  By
orthogonality, every choice of columns vanishes except columns of
degrees $0,\ldots,N-1$.  That surviving integral equals
$N!H_N^{(0)}=N!$.  The remaining minor is
$\det[p_{N+j}^{(0)}(z_i)]$.  Dividing by $\Delta(z)$ yields the identity
for the multiplier $\prod_\ell(z_\ell-x)$.

Finally put $z_\ell=-y_\ell$ and use
$p_{N+j}^{(0)}(-y)=(-1)^{N+j}R_{N+j}(y)$.  The signs from the multiplier,
column factors, and Vandermonde cancel, giving \eqref{eq:Christoffel}.
The numerator is alternating in the $y_i$, hence divisible by their
Vandermonde as a polynomial.  This proves the extension statement.
This is the classical fixed-size identity discussed in
\cite{Krattenthaler2023}; its inclusion here is a normalization proof,
not a novelty claim.
\end{proof}

Take $r=m+1$, let the first $m$ variables tend to zero, and put the last
one equal to $y=a/b\ne0$.  Dividing by the Vandermonde and taking the
limit replaces those first rows by Taylor coefficients.  Since
$[y^i]R_{N+j}(y)=\binom{N+j+i}{2i}$, we obtain:

\begin{proposition}[Confluent representation]\label{prop:confluent}
If $ab\ne0$, then
\begin{equation}\label{eq:smallD}
 \D{N}{m}(a,b)=b^Ny^{-m}
 \det\begin{pmatrix}
 \binom{N+j+i}{2i}\ \ (0\leq i<m,\ 0\leq j\leq m)\\[2pt]
 R_{N+j}(y)\ \ (0\leq j\leq m)
 \end{pmatrix},\qquad y=a/b.
\end{equation}
The upper block has $m$ rows and the lower block has one row.  For $m=0$
the formula simply reads $D_N^{(0)}=b^NR_N(a/b)$.
\end{proposition}

Define the $m\times m$ minors
\begin{equation}\label{eq:M}
 M_{m,\ell}(N)=
 \det_{\substack{0\leq i<m\\j\in J_\ell}}
       \binom{N+j+i}{2i},\qquad
 J_\ell=\{0,1,\ldots,m\}\setminus\{\ell\},
\end{equation}
where columns are in increasing order and $0\leq\ell\leq m$.
Expansion of \eqref{eq:smallD} along its last row gives
\begin{equation}\label{eq:minor-expansion}
 \D{N}{m}(a,b)=b^Ny^{-m}
 \sum_{\ell=0}^{m}(-1)^{m+\ell}M_{m,\ell}(N)R_{N+\ell}(y).
\end{equation}
Although a naive degree bound for $M_{m,\ell}$ is $m(m-1)=2d$, its
actual degree is only $d$.  This is the key point.

\section{The degree cancellation and its first correction}\label{sec:degree}

\begin{lemma}[Sharp minor degree]\label{lem:minor-degree}
For $0\leq\ell\leq m$, $M_{m,\ell}(N)$ is a polynomial of degree
$d=\binom m2$ with leading coefficient
\begin{equation}\label{eq:minor-leading}
 [N^d]M_{m,\ell}(N)=h_m\binom m\ell.
\end{equation}
For $d\geq1$, the next coefficient is
\begin{equation}\label{eq:minor-subleading}
 [N^{d-1}]M_{m,\ell}(N)
 =h_m\binom m\ell
   \left(\frac{d(m+2)}2-\frac{(m-1)\ell}{2}\right).
\end{equation}
\end{lemma}
\begin{proof}
Let $f_i(x)=\binom{x+i}{2i}$, so $f_i$ has degree $2i$ and leading
coefficient $1/(2i)!$.  The determinant $\det[f_i(x_j)]_{0\leq i,j<m}$
is alternating in its $m$ independent variables.  It is divisible by
$\Delta(x)=\prod_{p<q}(x_q-x_p)$, whose degree is $d$.  The determinant
has total degree at most $\sum_i2i=2d$; its quotient by $\Delta$ therefore
has total degree at most~$d$.  Substituting $x_j=N+j$ for
$j\in J_\ell$ makes $\Delta$ a nonzero constant.  Hence the degree in
$N$ is at most~$d$.

To obtain the leading coefficient, expand $f_i(N+j)$ in Taylor
powers of $j$.  The first possible nonzero determinant uses distinct
Taylor orders $0,1,\ldots,m-1$; these have the least possible sum $d$.
The coefficient of $N^d$ is
\begin{equation}\label{eq:Taylorleading}
 \frac{\det[\binom{2i}{r}]_{0\leq i,r<m}}{\prod_{i=0}^{m-1}(2i)!}
       \Delta(J_\ell)
 =\frac{2^d\Delta(J_\ell)}{\prod_{i=0}^{m-1}(2i)!}.
\end{equation}
To see the second equality, regard $\binom{2i}{r}$ as a polynomial
of degree $r$ in $i$ with leading coefficient $2^r/r!$ and evaluate
its determinant at $i=0,\ldots,m-1$.  The Vandermonde is
$\prod_{r=1}^{m-1}r!$, so these factorials cancel.
Removing $\ell$ from $0,\ldots,m$ also gives
\begin{equation}\label{eq:deletedVand}
 \Delta(J_\ell)=\binom m\ell\prod_{r=1}^{m-1}r!.
\end{equation}
Since
\[
 h_m=\frac{2^d\prod_{r=1}^{m-1}r!}{\prod_{i=0}^{m-1}(2i)!}
 \qquad(m\geq1),
\]
we obtain \eqref{eq:minor-leading}.  It is nonzero, proving the degree
assertion.  The case $m=0$ is the empty determinant and is immediate.

For the next coefficient, observe the symmetric factorization
\begin{equation}\label{eq:centered-binomial}
 f_i(x)=\frac1{(2i)!}\prod_{r=1}^{i}
          \left((x+\tfrac12)^2-(r-\tfrac12)^2\right).
\end{equation}
Thus $f_i(x)=(x+\tfrac12)^{2i}/(2i)!$ plus terms of degree at most
$2i-2$.  Replacing one or more rows by those lower-degree terms
reduces the total degree by at least two.  The same alternating
argument then bounds their degree after translation by $d-2$.
The leading two terms are consequently obtained from
\begin{equation}\label{eq:evenVand}
 \det\left[\frac{(N+j+\tfrac12)^{2i}}{(2i)!}\right]_
           {0\leq i<m,\ j\in J_\ell}
 =\frac{\displaystyle\prod_{\substack{p<q\\p,q\in J_\ell}}
        (q-p)(2N+p+q+1)}{\prod_{i=0}^{m-1}(2i)!}.
\end{equation}
The relative coefficient of $N^{d-1}$ is
\[
 \frac12\sum_{\substack{p<q\\p,q\in J_\ell}}(p+q+1)
 =\frac12\left((m-1)\left(\frac{m(m+1)}2-\ell\right)+d\right),
\]
which simplifies to \eqref{eq:minor-subleading}.  If $d=1$, the
putative remainder of degree $d-2=-1$ is identically zero.
\end{proof}

It is useful to collect the minors in a polynomial in an auxiliary
variable $r$:
\begin{equation}\label{eq:Q}
 Q_m(N;r)=\sum_{\ell=0}^{m}(-1)^{m+\ell}M_{m,\ell}(N)r^\ell.
\end{equation}
By the binomial theorem,
\begin{equation}\label{eq:Q-leading}
 [N^d]Q_m(N;r)=h_m(r-1)^m.
\end{equation}
When $r\ne1$ and $d\geq1$, the relative next coefficient is
\begin{equation}\label{eq:Q-subleading}
 \frac{[N^{d-1}]Q_m(N;r)}{[N^d]Q_m(N;r)}
 =\frac{d(m+2)}2-d\frac{r}{r-1}
 =d\left(\frac m2-\frac1{r-1}\right).
\end{equation}
Indeed, applying $r\,\dd/\dd r$ to $(r-1)^m$ sums the terms weighted
by $\ell$ in \eqref{eq:minor-subleading}.  The exact value $d$, rather
than an upper bound $2d$, is what makes the denominator exponent sharp.

\section{Two exponentials and the minimal recurrence}\label{sec:recurrence}

Suppose first that $b\ne0$ and $a(a+4b)\ne0$.  Set $y=a/b$ and choose
one ordering of the two roots
\begin{equation}\label{eq:rroots}
 r_\pm=\frac{y+2\pm\sqrt{y(y+4)}}2,
 \qquad r_+r_-=1.
\end{equation}
They are distinct and nonzero.  Solving the recurrence
\eqref{eq:Rrec} with its two initial values gives
\begin{equation}\label{eq:Rclosed}
 R_n(y)=\frac{(r_+-1)r_+^n-(r_--1)r_-^n}{r_+-r_-}.
\end{equation}
Write
\begin{equation}\label{eq:lambda}
 \lambda_\pm=br_\pm
    =\frac{a+2b\pm\sqrt{a(a+4b)}}2.
\end{equation}
Substitution in \eqref{eq:minor-expansion} gives the exact identity
\begin{equation}\label{eq:exactbranches}
 \D{N}{m}(a,b)=
 \frac{(r_+-1)Q_m(N;r_+)\lambda_+^N
       -(r_--1)Q_m(N;r_-)\lambda_-^N}
      {y^m(r_+-r_-)}.
\end{equation}
In particular, both polynomial coefficients in this two-exponential
representation have degree \emph{exactly} $d$: neither $r_+$ nor $r_-$
can equal $1$, since $y\ne0$.

\begin{lemma}[Minimal annihilator of an exponential polynomial]
\label{lem:minimal}
If $\lambda_1,\ldots,\lambda_s$ are distinct nonzero complex numbers
and the nonzero polynomials $A_i(N)$ have degrees $d_i$, then
$u_N=\sum_i A_i(N)\lambda_i^N$ has minimal constant-coefficient
annihilator
\begin{equation}\label{eq:generalminimal}
 \prod_{i=1}^s(E-\lambda_i)^{d_i+1},\qquad Eu_N=u_{N+1}.
\end{equation}
Its generating function has poles of orders $d_i+1$ at
$t=1/\lambda_i$.
\end{lemma}
\begin{proof}
On $A(N)\lambda^N$, the operator $E-\lambda$ is
$\lambda^{N+1}(A(N+1)-A(N))$; it lowers the polynomial degree by
exactly one until the result vanishes.  For $\mu\ne\lambda$,
$E-\mu$ preserves that degree, multiplying the leading coefficient
by $\lambda-\mu$.  The generalized eigenspaces for distinct eigenvalues
are independent: polynomial projectors from the Chinese remainder
identity isolate each one.  Thus every annihilator needs the stated
multiplicity at every eigenvalue.  This argument remains valid after
any finite shift in $N$.

Equivalently, applying $(z\,\dd/\dd z)^k$ to $1/(1-z)$ shows that
$\sum_{N\geq0}N^k z^N$ has a pole of exact order $k+1$ at $z=1$,
with nonzero leading pole coefficient $k!$.  Distinct pole locations
cannot cancel.  This proves the generating-function assertion too.
\end{proof}

Applying the lemma to \eqref{eq:exactbranches} proves the minimal
characteristic polynomial \eqref{eq:charpoly}.  It also proves that
$q^{d+1}F_m$ is a polynomial of degree at most $2d+1=2e-1$ for these
parameters.  Section~\ref{sec:numerator} sharpens its degree to $K$.

\subsection{Why the identity extends to all parameters}
Each $D_N^{(m)}$ is a polynomial in $a,b$ with integer coefficients.
The coefficients of $q^{d+1}F_m$ are therefore such polynomials too.
Every coefficient that vanishes for all $b\ne0$ and $a(a+4b)\ne0$
vanishes identically: a polynomial cannot be nonzero and vanish on
that nonempty open set.  Hence the uniform rational identity,
and its coefficient integrality, extend to every $a,b$, even where
its displayed denominator is not reduced.  This argument is
coefficientwise in a formal series and does not require a convergence
limit in $t$.

\subsection{Proof of the exceptional classification}
When $b=0$,
\begin{equation}\label{eq:bzero}
 \D{N}{m}(a,0)=a^N\HH{N}{m}.
\end{equation}
For $a\ne0$, the polynomial $H_N^{(m)}$ has degree $d$, so
Lemma~\ref{lem:minimal} gives denominator $(1-at)^{d+1}$.
When $a=0$ and $b\ne0$,
\begin{equation}\label{eq:azero}
 \D{N}{m}(0,b)=b^N\HH{N}{m+1},
\end{equation}
whose polynomial degree is $\binom{m+1}{2}$.

At the remaining collision $a=-4b$, with $b\ne0$, the recurrence gives
\begin{equation}\label{eq:Rminus4}
 R_n(-4)=(-1)^n(2n+1).
\end{equation}
Inserting this into \eqref{eq:minor-expansion} and tracking the signs
produces
\begin{align}
 \D{N}{m}(-4b,b)&=(-b)^N E_m(N),\label{eq:endpoint}\\
 E_m(N)&=4^{-m}\sum_{\ell=0}^{m}(2N+2\ell+1)M_{m,\ell}(N).
 \label{eq:E}
\end{align}
By \eqref{eq:minor-leading}, its degree is exactly $d+1$, with
\begin{equation}\label{eq:E-leading}
 [N^{d+1}]E_m(N)=4^{-m}\,2h_m\sum_{\ell=0}^m\binom m\ell
              =2^{1-m}h_m\ne0.
\end{equation}
Thus the reduced denominator is $(1+bt)^{d+2}$.
Finally, if $a=b=0$, every nonempty determinant vanishes and $F_m=1$.
Together with the nonexceptional result, this proves
Theorem~\ref{thm:exceptional}.

The two root-collision lines have different pole orders.  This is
not a contradiction: $x=0$ and $x=4$ are inequivalent endpoints of the
weight $x^{m-1/2}(4-x)^{1/2}$.  In particular, blindly substituting
$a=0$ into a generic minimal-order statement would be incorrect.

\section{The numerator degree and reciprocal symmetry}\label{sec:numerator}

A rational function can have a denominator of large degree and a
surprisingly short numerator.  Here the explanation is a string of
zero values in the natural bilateral extension of the determinant.

For the moment retain the nonexceptional hypotheses.  Define $D_N$ for
all $N\in\ZZ$ by the right side of \eqref{eq:smallD}.  The upper rows
are interpreted through their polynomial binomial coefficients; this
agrees with the Taylor coefficients of $R_n$ at zero, including
negative $n$, because
\[
 \binom{-n-1+i}{2i}=\binom{n+i}{2i}.
\]
It also agrees with the unique bilateral extension supplied by
\eqref{eq:exactbranches}.

\begin{lemma}[Reflection and a zero interval]\label{lem:bilateral}
Let $\sigma_m=(-1)^{\binom{m+1}{2}}$.  Then
\begin{align}
 D_{-N-m-1}&=\sigma_m b^{-2N-m-1}D_N\qquad(N\in\ZZ),
       \label{eq:Dreflect}\\
 D_{-1}=D_{-2}=\cdots=D_{-m}&=0.
       \label{eq:Dzeros}
\end{align}
The second statement is empty for $m=0$.
\end{lemma}
\begin{proof}
The index change $N\mapsto-N-m-1$, followed by
$R_{-n-1}=R_n$, reverses all $m+1$ columns of the determinant.
The sign of that reversal is $\sigma_m$.  Comparing its prefactor
$b^N$ on the two sides gives \eqref{eq:Dreflect}.
If $N=-s$ with $1\leq s\leq m$, columns $j=s-1$ and $j=s$ correspond
to $R_{-1}$ and $R_0$, respectively, which are the same polynomial.
Their values and all needed Taylor coefficients coincide, so the
determinant vanishes.
\end{proof}

\begin{lemma}[The expansion at infinity]\label{lem:infinity}
For a sequence of exponential-polynomial form with nonzero characteristic
roots and its bilateral extension, its proper rational generating
function satisfies
\begin{equation}\label{eq:infinity}
 F(t)=-\sum_{r\geq1}D_{-r}t^{-r}
\end{equation}
as a formal Laurent expansion at infinity.
\end{lemma}
\begin{proof}
For $D_N=\lambda^N$ the assertion is simply
$(1-\lambda t)^{-1}=-\sum_{r\geq1}\lambda^{-r}t^{-r}$.
Apply powers of $t\,\dd/\dd t$ to obtain the assertion for
$D_N=N^k\lambda^N$, and then use linearity.
\end{proof}

Combining \eqref{eq:Dzeros}, \eqref{eq:Dreflect}, and $D_0=1$ yields
\begin{equation}\label{eq:leadinginfinity}
 F_m(t)=-\sigma_m b^{-m-1}t^{-m-1}+O(t^{-m-2})
 \qquad(t\longrightarrow\infty).
\end{equation}
Since $q^e$ has degree $2e$ and leading coefficient $b^{2e}$, the
numerator has degree
\begin{equation}\label{eq:degreeK}
 2e-m-1=m(m-1)+2-m-1=(m-1)^2=K
\end{equation}
and leading coefficient $-\sigma_m b^{2e-m-1}=s_m b^K$.
The identity $s_m=-\sigma_m$ follows because the difference of the two
exponents $\binom{m+1}{2}$ and $(m-1)(m-2)/2$ is $2m-1$, an odd integer.
This calculation proves the sharp degree assertion, including $m=0$.

The full reflection identity yields more than the leading term.
Putting $1/(b^2t)$ into \eqref{eq:infinity} and reindexing with
\eqref{eq:Dreflect}, we find
\begin{equation}\label{eq:Freflect}
 F_m\!\left(\frac1{b^2t}\right)
 =s_m b^{m+1}t^{m+1}F_m(t).
\end{equation}
But $q(1/(b^2t))=q(t)/(b^2t^2)$.  Substitution into
\eqref{eq:Freflect} gives exactly \eqref{eq:reciprocity}.
After denominators in $b$ are cleared, each of these numerator identities
is polynomial in $a,b$.  Thus they extend to the exceptional parameters;
for $b\ne0$ the leading coefficient remains nonzero.  This completes the
proof of Theorem~\ref{thm:main}.

\subsection{An explicit finite prescription for the numerator}
No uncomputed infinite series is needed to determine $P_m$.  Expand
\begin{equation}\label{eq:qcoeffs}
 q(t)^e=\sum_{j=0}^{2e}q_jt^j.
\end{equation}
For example, its coefficients have the finite expression
\begin{equation}\label{eq:qexplicit}
 q_j=\sum_{v=\max(0,j-e)}^{\lfloor j/2\rfloor}
 \frac{e!}{(j-2v)!v!(e-j+v)!}
       (-(a+2b))^{j-2v}b^{2v}.
\end{equation}
Then
\begin{equation}\label{eq:Pcompute}
 P_m(t)=\sum_{\ell=0}^{K}
 \left(\sum_{j=0}^{\min(\ell,2e)}q_j\D{\ell-j}{m}(a,b)\right)t^\ell,
\end{equation}
where every determinant in the right side is given by
Theorem~\ref{thm:coeff}.  Equations~\eqref{eq:qexplicit} and
\eqref{eq:Pcompute} therefore constitute a finite sum/product formula
for every numerator coefficient.  All coefficients of $q^eF_m$ after
degree $K$ vanish.

The usual recurrence is
\begin{equation}\label{eq:recurrence}
 \sum_{j=0}^{2e}q_j\D{N-j}{m}(a,b)=0\qquad(N\geq2e).
\end{equation}
In fact the truncated convolution already vanishes for $N>K$ when
negative-index terms are taken as zero for the purpose of multiplying
ordinary power series.  Those artificial zeros should not be confused
with the distinct bilateral extension used in the proof above.

\section{Exact exponential expansions and asymptotics}\label{sec:asymptotics}

Equation~\eqref{eq:exactbranches} supplies every coefficient of the two
polynomials in Theorem~\ref{thm:asymp}.  Explicitly,
\begin{equation}\label{eq:Bpolys}
 B_+(N)=\frac{(r_+-1)Q_m(N;r_+)}{y^m(r_+-r_-)},\qquad
 B_-(N)=-\frac{(r_--1)Q_m(N;r_-)}{y^m(r_+-r_-)}.
\end{equation}
These formulas are valid for arbitrary nonexceptional complex parameters;
the assumption $a,b>0$ is only needed to select a strictly dominant real
exponential for the asymptotic statement.

For $a,b>0$, write $r=r_+$, so $r_-=r^{-1}$ and
\begin{equation}\label{eq:yparam}
 y=r+r^{-1}-2=\frac{(r-1)^2}{r}.
\end{equation}
Using \eqref{eq:Q-leading}, the coefficient of $N^d$ in $B_+$ is
\[
 \frac{h_m(r-1)^{m+1}}{y^m(r-r^{-1})}
 =h_m\frac{r^{m+1}}{(r-1)^m(r+1)}=A_m.
\]
Its relative next coefficient is exactly \eqref{eq:Q-subleading}.
Also $\lambda_-/\lambda_+=r^{-2}$.  Division of
\eqref{eq:exactbranches} by $A_mN^d\lambda_+^N$ proves
\eqref{eq:asymp}.

One may write the complete answer in the finite form
\begin{equation}\label{eq:fullfinite}
 \D{N}{m}=\lambda_+^N\sum_{j=0}^{d}\beta_j^+N^{d-j}
          +\lambda_-^N\sum_{j=0}^{d}\beta_j^-N^{d-j},
\end{equation}
where
\begin{equation}\label{eq:allbetas}
 \beta_j^+=\frac{r_+-1}{y^m(r_+-r_-)}[N^{d-j}]Q_m(N;r_+),\qquad
 \beta_j^-=-\frac{r_--1}{y^m(r_+-r_-)}[N^{d-j}]Q_m(N;r_-).
\end{equation}
Thus the ``full asymptotic expansion'' is particularly strong here:
it is an exact expression, with no uncomputed remainder once both
finite sectors are retained.  No assertion of asymptotic convergence
or summability is needed.  The leading coefficient ratio of the two
sectors is
\begin{equation}\label{eq:branchratio}
 \frac{\beta_0^-}{\beta_0^+}=(-1)^m r^{-(m+1)}.
\end{equation}

These estimates are for fixed $m,a,b$.  They are not claimed to be
uniform when $a\to0$, when the two roots coalesce, or when $m$ grows
with $N$.  In such regimes the exact finite formula remains valid
away from the collision itself, but a different uniform asymptotic
analysis would be required.  At a collision, use the exact exceptional
formulas from Section~\ref{sec:recurrence} instead.

\section{Worked examples and integer sequences}\label{sec:examples}

\subsection{Small shifts}
The first uniform generating functions are
\begin{align}
 F_0(t)&=\frac{1-bt}{q(t)},&
 F_1(t)&=\frac1{q(t)},\label{eq:low01}\\
 F_2(t)&=\frac{1+bt}{q(t)^2},&
 F_3(t)&=\frac{(1-b^2t^2)(1+(a+6b)t+b^2t^2)}{q(t)^4}.
 \label{eq:low23}
\end{align}
They can be obtained directly from \eqref{eq:Pcompute}.  Notice that
$m=0$ is a genuine boundary case: its numerator has degree $1$, not $0$.
For $m=3$, the numerator is anti-reciprocal; for $m=2$, it is reciprocal.
This is the pattern described uniformly by $s_m$.

For $m=2$, put $A=a+2b$.  The recurrence becomes
\begin{equation}\label{eq:m2rec}
 D_N=2A D_{N-1}-(A^2+2b^2)D_{N-2}
         +2Ab^2D_{N-3}-b^4D_{N-4}\qquad(N\geq4).
\end{equation}
The initial polynomials are
\begin{align*}
 D_0&=1,\\
 D_1&=2a+5b,\\
 D_2&=3a^2+14ab+14b^2,\\
 D_3&=4a^3+27a^2b+54ab^2+30b^3.
\end{align*}
The recurrence is minimal outside the exceptional lines; displaying
only a valid order-four relation would not by itself prove that fact.

\subsection{The specialization \texorpdfstring{$a=b=1$}{a=b=1}}
The first sequences, all beginning with the empty determinant, are
\begin{center}
\renewcommand{\arraystretch}{1.2}
\begin{tabular}{@{}r rrrrrrr@{}}
\toprule
$m\backslash N$ & 0 & 1 & 2 & 3 & 4 & 5 & 6\\
\midrule
0 & 1 & 2 & 5 & 13 & 34 & 89 & 233\\
1 & 1 & 3 & 8 & 21 & 55 & 144 & 377\\
2 & 1 & 7 & 31 & 115 & 390 & 1254 & 3893\\
3 & 1 & 19 & 170 & 1075 & 5580 & 25529 & 107036\\
\bottomrule
\end{tabular}
\end{center}
The first two are $F_{2N+1}$ and $F_{2N+2}$, respectively, as follows
from their common recurrence $u_{N+2}=3u_{N+1}-u_N$ and initial values.
They connect the family to the odd and even Fibonacci subsequences in
OEIS A001519 and A001906~\cite{OEISfib}; Catalan numbers themselves are
A000108~\cite{OEIScat}.  The higher rows and their coefficient triangles
are natural OEIS-style objects, but the archive does not assert that
any particular row is absent from OEIS.

For every fixed shift, the dominant exponential is
\[
 \lambda_+=\frac{3+\sqrt5}{2},
\]
while its polynomial prefactor has degree $\binom m2$.  Thus increasing
the shift changes the power of $N$, not the exponential growth constant.
For $m=2$, the first relative asymptotic correction equals
$1-1/(r-1)=(3-\sqrt5)/2$ with $r=(3+\sqrt5)/2$.

\subsection{Rational characteristic roots: \texorpdfstring{$a=1,b=2$}{a=1,b=2}}
This specialization gives $r_+=2$, $r_-=1/2$ and
$\lambda_+=4$, $\lambda_-=1$.  Therefore the exact formulas have
rational polynomial coefficients and no radicals:
\begin{align}
 \D{N}{0}(1,2)&=\frac{2\cdot4^N+1}{3},\label{eq:example0}\\
 \D{N}{1}(1,2)&=\frac{4\cdot4^N-1}{3},\label{eq:example1}\\
 \D{N}{2}(1,2)&=\frac{8N\,4^N+N+3}{3},\label{eq:example2}\\
 \D{N}{3}(1,2)&=
 \frac{(32N^3+48N^2+16N+96)4^N
       -(2N^3+21N^2+73N+78)}{18}.
 \label{eq:example3}
\end{align}
For example, the $m=2$ sequence starts
$1,12,87,514,2733,13656,\ldots$.  Its dominant polynomial is
$(8/3)N$ with no constant term.  This vanishing agrees with the
first-correction formula $d(m/2-1/(r-1))=1(1-1)=0$; the entire
correction is the smaller exponential sector $(N+3)/3$.
For $m=3$, the relative coefficient of $N^{-1}$ in the dominant sector
is $3/2$, again exactly as predicted.

\subsection{An endpoint example}
For $m=2$, \eqref{eq:E} simplifies to
\[
 E_2(N)=\frac{(N+1)(N+2)}2.
\]
Hence
\begin{equation}\label{eq:endpointexample}
 \D{N}{2}(-4b,b)=(-b)^N\frac{(N+1)(N+2)}2,
 \qquad F_2(t;-4b,b)=\frac1{(1+bt)^3}.
\end{equation}
This also follows by cancelling $1+bt$ from
$(1+bt)/(1+bt)^4$, the uniform form of $F_2$.
At the other collision, $a=0$, one instead obtains
$F_2(t;0,b)=(1+bt)/(1-bt)^4$.  These explicit cases illustrate why the
reduced denominator needs a separate classification.

\section{Algorithms, exact experiments, and reproducibility}\label{sec:code}

\subsection{Implementation choices}
The file \texttt{code/catalan\_hankel.py} implements the coefficient
formula, pure shifted product, homogeneous polynomial evaluation,
fixed-size determinant, and uniform and reduced denominators.  All
numerical parameters in its public computation interface are integers;
internal fractions use Python's exact \texttt{Fraction}.  The
mathematical theorems have the broader complex-parameter scope stated
above.  No floating-point arithmetic enters the tests.

For a coefficient row, first compute the triangular product
$H_N^{(m+1)}$ and then apply \eqref{eq:coeff-update} successively.
The number of arithmetic operations is $O(m^2+N)$; this is an
\emph{arithmetic-operation count}, not a bit-complexity claim.
Every division in the update is checked for exactness.  Homogeneous
Horner evaluation avoids division by $a$ or $b$, so zero parameters
are handled without special numerical limits.

The independent determinant path uses a fraction-free Bareiss algorithm
on the original matrix of Catalan moments, including pivoting when a
diagonal pivot vanishes.  This matters for signed parameters: an
implementation tested only on positive definite examples can conceal
pivot failures.  The source-indexing formula is implemented separately,
not obtained by calling the corrected coefficient routine.

For many consecutive values at fixed $m$, it is also possible to use
\eqref{eq:recurrence} after producing the first $2e$ values.  For very
large isolated indices, powering the companion matrix or reducing
$X^N$ modulo \eqref{eq:charpoly} provides standard recurrence-based
methods.  These methods are consequences of the theorem; the supplied
reference implementation prioritizes transparent exact verification
rather than optimized large-integer throughput.

\subsection{Verification that was actually run}
The included \texttt{data/verification.json} records \texttt{PASS} for
all checks in the following table.  The report is generated by the test
program, not manually written.
\begin{center}
\small
\renewcommand{\arraystretch}{1.17}
\begin{tabular}{@{}p{0.68\textwidth}r@{}}
\toprule
Check & Count\\
\midrule
Direct determinant versus coefficient evaluation & 1,183\\
Additional direct signed/degenerate parameter evaluations & 810\\
Printed source formula versus corrected shift & 990\\
Confluent fixed-size determinant evaluations & 364\\
Minor polynomials exactly interpolated and checked & 36\\
Generating-function/recurrence coefficient identities & 2,816\\
Numerator reciprocal coefficient identities & 604\\
Nonexceptional numerator--denominator gcd checks & 44\\
Exceptional reduced-denominator checks & 44\\
Exact two-exponential evaluations & 168\\
Strict normalized log-concavity inequalities & 5,655\\
\bottomrule
\end{tabular}
\end{center}

The first 1,183 evaluations cover every pair $0\leq N,m\leq12$,
using $a=0,1,\ldots,N$ and $b=1$.  At each fixed $(N,m)$, both sides
are polynomials of degree at most $N$ in $a$.  Equality at $N+1$
distinct points certifies that particular polynomial identity exactly.
Thus this part certifies 169 full coefficient polynomials, not merely
169 evaluations.  It still does not prove the result for untested
values of $N$ or $m$; the proof in Section~\ref{sec:coeffproof} does that.

For $0\leq m\leq7$, the minor tests interpolate using $2d+2$ points,
more than enough for the elementary degree bound $2d$.  They verify
the drop to degree $d$, both leading coefficients from
Lemma~\ref{lem:minor-degree}, and additional negative-index values.
The recurrence checks include $0\leq m\leq10$ and four nonexceptional
parameter pairs, while separate tests examine every exceptional class.
The exact two-exponential tests use $(a,b)=(1,2)$, $m\leq7$, and
$0\leq N\leq20$.  These tests check different parts of the derivation,
but are not represented as formal proof certificates for arbitrary
parameters.

\subsection{How to reproduce the package}
From the archive's top directory, run:
\begin{lstlisting}
python3 code/verify.py --out data
python3 code/catalan_hankel.py 8 3 --a 2 --b 3 --direct
latexmk -pdf -interaction=nonstopmode -halt-on-error article.tex
\end{lstlisting}
Use Python 3.10 or later and do not pass \texttt{-O}, because the
verification suite deliberately uses assertions.  The script writes
coefficient rows, generating functions, minor polynomials, exact
exponential-sector polynomials, and integer sequence tables in JSON or
CSV.  Rational polynomial coefficients are stored as strings such as
\texttt{"8/3"}, in ascending power order, to avoid any loss of exactness
when imported into other systems.  The included README documents every
file and its indexing convention.

\section{What has been resolved, and what has not}\label{sec:status}

The all-shift denominator formula is proved here, with a stronger
minimality assertion and a complete description of the degenerate
parameters.  The coefficient expression is likewise proved for all
$N,m$; comparison with the exact formula printed in the target source
requires the explicit indexing correction in the next appendix.
The reciprocal numerator and finite exponential expansions are proved
from the same structural determinant calculation.

These are substantive mathematical conclusions about the selected
conjectural family.  They are not evidence that the Christoffel identity,
Jacobi representation, Catalan shifted product, or general rationality
are new.  A literature search was performed against the source title,
arXiv identifier, denominator phrase, and related recurrence work.
The cited version still labels the sharp exponent conjectural, and the
retrieved sources did not supply an explicit statement of this full
sharp classification.  That is a bounded search result, not a proof
that no earlier equivalent theorem exists.  Before a priority claim
or submission, the exact formulas should be checked against the wider
orthogonal-polynomial and determinant literature and by a human referee.

Several natural extensions remain separate from the theorem.  One may
seek a direct combinatorial interpretation of the individual positive
integers $c_{N,m,k}$, or a product formula for the minors
$M_{m,\ell}(N)$ beyond their first two coefficients.  Multiple distinct
linear factors in the multiplier lead to more exponential sectors and
new collision patterns.  A joint limit with growing $m$ or a parameter
approaching $a=0$ would require uniform, rather than fixed-parameter,
asymptotics.  None of these extensions is assumed in the present proofs.

\appendix
\section{The precise correction to the printed conjecture}\label{app:barry}

This appendix distinguishes two kinds of index.  The source uses a
size-$(n+1)$ determinant.  We write $N=n+1$.  To avoid silently
identifying incompatible shift conventions, write $s$ for the source's
symbol in its displayed coefficient formula, and reserve $m$ for
our actual Catalan shift.

For $s\geq2$ and $0\leq k\leq n+1$, the coefficient expression printed
in Conjecture~2 of~\cite{Barry2020} can be written as
\begin{align}
 T_{n,k,s}
 &=\frac{C_s}{2s-2}
   \binom{s+k-2}{s-2}\binom{n+k+2s-2}{2k+2s-3}\nonumber\\[-1pt]
 &\quad\times
  \prod_{j=0}^{\lfloor(s-1)/2\rfloor-1}
       \frac{\binom{2n+2s-2j-1}{2s-4j-5}}
            {\binom{2s-2j-1}{2s-4j-5}}
  \prod_{j=0}^{s-3}\frac{2s-j-2}{n+k+2s-j-2}.
 \label{eq:BarryT}
\end{align}
Only the natural, nonambiguous range $s\geq2$ is needed here.  Empty
products have value~$1$.

\begin{proposition}[Correct shift]\label{prop:Barry}
For $N\geq1$, $s\geq2$, and $0\leq k\leq N$,
\begin{equation}\label{eq:Barrycorrect}
 T_{N-1,k,s}=c_{N,s-1,k}.
\end{equation}
Consequently, the polynomial using $T_{N-1,k,s}$ is
$D_N^{(s-1)}(a,b)$, not $D_N^{(s)}(a,b)$.  Equivalently, to use that
printed expression for an actual shift $m\geq1$, replace its shift
symbol by $s=m+1$.
\end{proposition}
\begin{proof}
Taking a ratio of successive terms in \eqref{eq:BarryT} and cancelling
factorials gives
\begin{equation}\label{eq:Tkratio}
 \frac{T_{n,k+1,s}}{T_{n,k,s}}
 =\frac{(n+1-k)(n+k+s+1)}{2(k+1)(2k+2s-1)}.
\end{equation}
This agrees with \eqref{eq:coeff-update} after $N=n+1$ and $m=s-1$.
It remains to match the initial coefficient.
Putting $n=N-1$, $k=0$, and simplifying the factorials gives
\begin{equation}\label{eq:TB0}
 T_{N-1,0,s}=C_s\frac{(N+s-1)!}{N!s!}
 \prod_{j=0}^{p-1}
 \frac{\binom{2N+2s-2j-3}{2s-4j-5}}
      {\binom{2s-2j-1}{2s-4j-5}},
 \qquad p=\left\lfloor\frac{s-1}{2}\right\rfloor.
\end{equation}
At $N=1$ this is $C_s=H_1^{(s)}$.  Its ratio at successive $N$ is
\begin{align}
 \frac{T_{N,0,s}}{T_{N-1,0,s}}
 &=\frac{N+s}{N+1}
   \prod_{j=0}^{p-1}
   \frac{(2N+2s-2j-1)(2N+2s-2j-2)}
        {(2N+2j+4)(2N+2j+3)}\nonumber\\
 &=\frac{(2N+2s)!(2N+1)!}
         {(2N+s)!(2N+s+1)!}.
 \label{eq:TNr}
\end{align}
In the product, the numerator factors form a consecutive interval of
integers, as do the denominator factors; the last equality holds for
either parity of $s$.  The shifted-product formula gives the identical
ratio
\[
 \frac{H_{N+1}^{(s)}}{H_N^{(s)}}
 =\prod_{i=1}^{s-1}
     \frac{(2N+2i+1)(2N+2i+2)}{(2N+i+1)(2N+i+2)}
 =\frac{(2N+2s)!(2N+1)!}{(2N+s)!(2N+s+1)!}.
\]
Induction in $N$ therefore proves $T_{N-1,0,s}=H_N^{(s)}$.
The ratio \eqref{eq:Tkratio} then proves \eqref{eq:Barrycorrect}
for every $k$.
\end{proof}

The discrepancy is visible in the smallest nontrivial instance.
With $N=1$ and $s=2$, the printed expression is
\[
 T_{0,0,2}b+T_{0,1,2}a=2b+a,
\]
whereas the determinant with the stated shift $s=2$ is
\[
 D_1^{(2)}(a,b)=aC_2+bC_3=2a+5b.
\]
At $a=b=1$, these are $3$ and $7$.  The expression $a+2b$ is exactly
$D_1^{(1)}(a,b)$, so the correction is systematic rather than an isolated
numerical exception.  The original PDF page was inspected as well as
its HTML rendering; this is not an artifact of an HTML conversion.
The denominator-exponent statement is separate and uses the actual
shift $m$ correctly.

\section{Compact reference and archive conventions}\label{app:reference}

\begin{center}
\renewcommand{\arraystretch}{1.25}
\begin{tabular}{@{}p{0.22\textwidth}p{0.72\textwidth}@{}}
\toprule
Symbol & Meaning\\
\midrule
$N$ & Determinant size, with $D_0=1$\\
$m$ & Actual Catalan shift in $aC_{i+j+m}+bC_{i+j+m+1}$\\
$H_N^{(m)}$ & Pure shifted Catalan determinant\\
$d,e,K$ & $\binom m2$, $\binom m2+1$, and $(m-1)^2$\\
$P_m$ & Uniform numerator, not automatically reduced at exceptions\\
$q$ & $1-(a+2b)t+b^2t^2$\\
$M_{m,\ell}$ & Minor with column $\ell$ deleted from the binomial block\\
$Q_m(N;r)$ & Signed generating polynomial of those minors\\
$r_\pm$ & Roots of $r^2-(a/b+2)r+1$\\
$\lambda_\pm$ & $br_\pm$, the characteristic roots in determinant size\\
$s$ in Appendix A & Source shift symbol, equal to actual shift plus one\\
\bottomrule
\end{tabular}
\end{center}

The JSON files use ascending polynomial powers.  In
\texttt{coefficient\_rows.json}, position $k$ is the coefficient of
$a^kb^{N-k}$.  In \texttt{minor\_polynomials.json}, position $j$ is the
coefficient of $N^j$.  In \texttt{two\_branch\_polynomials.json}, the
arrays \texttt{P\_plus} and \texttt{P\_minus} describe the polynomial
multipliers of $4^N$ and $1^N$ for $(a,b)=(1,2)$; they are not the
ordinary-generating-function numerator $P_m$.

The source audit records exact URLs, source versions, and search
limitations.  The tests establish reproducibility of the computation,
not priority or formal verification.  The archive intentionally contains
no third-party article PDFs, no font files, no compiler auxiliaries,
and no Python bytecode.

\clearpage
\begin{thebibliography}{9}
\addcontentsline{toc}{section}{References}

\bibitem{Barry2020}
P.~Barry,
\emph{Notes on the Hankel transform of linear combinations of consecutive
pairs of Catalan numbers}, arXiv:2011.10827v1 (21 November 2020), 29 pp.
The denominator conjecture and Conjecture~2 are on p.~4.
\url{https://arxiv.org/abs/2011.10827}.

\bibitem{Dougherty2011}
M.~Dougherty, C.~French, B.~Saderholm, and W.~Qian,
Hankel transforms of linear combinations of Catalan numbers,
\emph{Journal of Integer Sequences} \textbf{14} (2011), Article 11.5.1.
\url{https://cs.uwaterloo.ca/journals/JIS/VOL14/French/french2.html}.

\bibitem{Krattenthaler2023}
C.~Krattenthaler,
Hankel determinants of linear combinations of moments of orthogonal
polynomials, II,
\emph{The Ramanujan Journal} \textbf{61} (2023), 597--627;
published online 22 November 2021.
\url{https://doi.org/10.1007/s11139-021-00514-8}.
See the determinant formula, the confluence discussion, and Section~8
for existing general recurrence theory and its earlier sources.

\bibitem{DLMF}
NIST Digital Library of Mathematical Functions,
Section~18.5, Eq.~18.5.7, Jacobi polynomial hypergeometric representation.
\url{https://dlmf.nist.gov/18.5.E7}.

\bibitem{OEIScat}
OEIS Foundation Inc., \emph{The On-Line Encyclopedia of Integer Sequences},
A000108 (Catalan numbers).
\url{https://oeis.org/A000108}.

\bibitem{OEISfib}
OEIS Foundation Inc., \emph{The On-Line Encyclopedia of Integer Sequences},
A001519 and A001906 (Fibonacci subsequences).
\url{https://oeis.org/A001519}; \url{https://oeis.org/A001906}.

\end{thebibliography}
\noindent\small\textit{Online sources checked on 19 September 2026.
All theorem numbering in the body refers to the present manuscript,
unless a source is explicitly named.}
\end{document}
